\subsection{Results on Generalizability and Robustness to Adversarial Attacks}
\tit{Current Deepfake Detectors do not Generalize} Most existing deepfake detectors rely on binary classification to distinguish between real and fake images. However, as shown in ~\Cref{fig:c2_clip}, this approach often leads to a strong bias toward classifying out-of-domain inputs as real. In particular, the C2-CLIP detector correctly classifies only 51 fake images, while misclassifying 949 as real, despite correctly identifying most authentic inputs. This imbalance highlights a strong bias toward high recall on real images at the cost of failing to flag synthetic content. A similar failure pattern is evident under adversarial conditions. As shown in~\Cref{fig:d3_pgd} and~\Cref{tab:attack_results}, all evaluated detectors collapse completely under imperceptible PGD~\cite{pgd} perturbations, with several models including UFD~\cite{ojha2023towards}, FreqNet~\cite{Tan_Zhao_Wei_Gu_Liu_Wei_2024}, and FatFormer~\cite{fatformer} dropping to 0\% accuracy post-attack. Even D3~\cite{d3}, the strongest baseline, falls from 83.90\% to just 1.75\% accuracy. These results underscore a general and critical weakness: traditional detectors exhibit a strong confirmation bias toward authenticity and offer no graceful degradation under adversarial manipulation or distributional shift. Refer to supplementary for additional results.
